\documentclass[12pt]{article}
\usepackage[utf8]{inputenc}
\usepackage[T1]{fontenc}
\usepackage{hyperref}
\title{Awa Language Resources for K-12 Curriculum Development}
\author{Compiled by Claude Sonnet 5 (claude-sonnet-5) for Andr\'as Kornai}
\date{}
\begin{document}
\maketitle
\section{Introduction}
Awa (ISO 639-3 vwa; Glottolog awac1238) is a Waic (Palaungic, Austroasiatic) variety spoken in Yunnan, China. Ethnologue-style sources group Parauk, Vo and Awa as varieties of Wa and estimate about 100,000 Awa speakers. Jackson Sun argues the Chinese ``Awa'' label lumps distinct languages such as Va and Lavia, so classification is disputed. The standard Wa script used in China is based on the Aishuai variety of Cangyuan County.
\section{Primary Educational Resources}
No Awa-specific primers or textbooks were found. Any teaching materials found relate to standard Wa, not to Awa.
\section{Linguistic and Reference Materials}
General overview of the Wa language: \url{https://en.wikipedia.org/wiki/Wa_language}. Sun's paper on the Awa grouping and the Va language: \url{https://www.jbe-platform.com/content/journals/10.1075/ltba.18010.sun}
\section{Religious and Community Resources}
None found in search results.
\section{Digital and Audio Resources}
None specific to Awa found in search results.
\section{Curriculum Development Context}
Little exists for Awa itself. A curriculum would probably have to draw on standard Wa materials and new documentation of Awa. The boundaries between Awa, Va and other varieties should be settled with the community first.
\end{document}